\documentclass[10pt,a4paper]{amsart}
\usepackage[margin=19mm]{geometry}
\usepackage{amsmath,amssymb,mathtools}
\usepackage[scr=rsfs,cal=euler]{mathalfa}
\usepackage{xfrac}
\usepackage[unicode,hidelinks,bookmarks=false]{hyperref}
\newcommand{\ie}{i.e.\ }
\newcommand{\cf}{cf.\ }
\newcommand{\resp}{respectively\ }
\newcommand{\Subsection}{\subsection}
\providecommand{\nobreakdash}{\nobreak\hbox{-}}
\setlength{\emergencystretch}{2em}
\multlinegap0pt
\usepackage{bm}
\usepackage{bbm}
\usepackage{dsfont}
\usepackage{xspace}
\usepackage{cancel}
\usepackage{mathtools}
\usepackage{amsthm}
\makeatletter
\@ifundefined{theorem}{\newtheorem{theorem}{Theorem}}{}
\@ifundefined{lemma}{\newtheorem{lemma}[theorem]{Lemma}}{}
\makeatother
\DeclareMathOperator{\Prob}{Prob}
\newcommand{\cO}{\mathcal{O}}
\newcommand{\cT}{\mathcal{T}}
\newcommand{\ra}{\rightarrow}
\title[A Szemer\'edi type theorem for sets of positive density in a]{A Szemer\'edi type theorem for sets of positive density in approximate lattices}
\author{Michael Bj\"orklund, Alexander Fish}
\date{Source: arXiv:2402.17158v3 (CC BY 4.0)}
\begin{document}
\maketitle

\noindent\emph{Source and licence.} A Szemer\'edi type theorem for sets of positive density in approximate lattices. Version arXiv:2402.17158v3 (CC BY 4.0), distributed under the Creative Commons Attribution 4.0 International licence, as recorded in the arXiv licence metadata for this exact version and shown on the abstract page https://arxiv.org/abs/2402.17158v3. Full rights notice: licenses/arxiv-2402.17158-CC-BY-4.0.txt.\par\smallskip

\noindent\textit{Editorial scope.} Counted unit: the Lemma whose source label is \texttt{lemma3} in the article (source0.tex lines 1242--1252, source label \texttt{lemma3}), delivered together with the article's own complete original proof of it (source0.tex lines 1254--1306, one \texttt{proof} environment of 53 lines). No mathematics is written, repaired or re-derived: statement and proof are the article's own text line for line. Required context retained uncounted: Lemma\_g1gr (lines 785--797); Lemma\_TransSep (lines 862--870). Every definition, display and label of those lines is retained unchanged. Omitted from this package and not redistributed: 1--784, 798--861, 871--1241, 1253--1253, 1307--1973. Nothing inside a retained excerpt is omitted or altered: every retained body is delivered exactly as the article's lines are written, so no omission receipt of any kind accompanies this package. Citations: the retained excerpts carry no citation command at all, so no bibliography entry of the article is transcribed and no reference is redistributed. Reference adaptation: none was needed and none was made; every reference inside the delivered unit resolves inside it, so no unresolved reference or citation is produced. Printed slips kept literal: no printed word, symbol, hypothesis or number of the counted statement or of its proof was altered. Layout and class: the article's own class is \texttt{amsart} with packages including mdframed, setspace, mathrsfs, comment, inputenc, hyperref, bold-extra, etoolbox, amsmath, enumerate, amssymb, amscd, amsthm, amsfonts; the wrapper is the lane's pinned \texttt{amsart} editorial wrapper at 10pt on A4, which restates the theorem-like environments the retained text needs and adds the article's own macro definitions that the retained text needs (\texttt{\textbackslash Prob}, \texttt{\textbackslash cO}, \texttt{\textbackslash cT}, \texttt{\textbackslash ra}), with extra wrapper packages (bm, bbm, dsfont, xspace, cancel, mathtools). The delivered numbering is the unit's own. Assets: no retained span contains a figure environment, a graphics-inclusion command, a PDF-page inclusion, a table or a TikZ picture; no image file is copied into this package, the package declares no asset dependency and no image file is redistributed with it. Licence: version arXiv:2402.17158v3 of this article is distributed under the Creative Commons Attribution 4.0 International licence, as recorded in the arXiv licence metadata for this exact version and shown on the abstract page https://arxiv.org/abs/2402.17158v3. A full rights notice is delivered at licenses/arxiv-2402.17158-CC-BY-4.0.txt. Characters: every retained excerpt is pure ASCII, so the delivered engine, XeLaTeX with its default Latin Modern text font, reports no missing character.
\begin{lemma}
\label{Lemma_g1gr}
Let $\Delta \subset G$ and suppose that there is an open identity neighbourhood $W \subset U$ such that $(\Xi_Y - \Delta) \cap W = \{0\}$. For every Borel set $V_o$ in $G$ 
such that $V_o-V_o \subset W$, and for every Borel set $B \subset Y$, 
\[
V_o.\Big(B \cap \bigcap_{k=1}^r (-g_k).B\Big) = B_{V_o} \cap \Big( \bigcap_{k=1}^r (-g_k).B_{V_o}\Big), \quad \textrm{for all $g_1,\ldots,g_r \in \Delta$},
\]
where $B_{V_o} = V_o.B$. In particular, for every $\mu \in \Prob_G(X)$, 
\[
\mu_Y\Big( B \cap \bigcap_{k=1}^r (-g_k).B\Big)
= \frac{1}{m_G(V_o)} \cdot \mu\Big(B_{V_o} \cap \Big( \bigcap_{k=1}^r (-g_k).B_{V_o}\Big)\Big).
\]
\end{lemma}

\begin{lemma}
\label{Lemma_TransSep}
Let $P \in \Omega_{P_o}$. Then,  
\[
P-P \subset \overline{P_o - P_o} \quad \textrm{for all $P \in \Omega^{\times}_{P_o}$},
\] 
and $\cT_{P_o}$ is a $U$-separated cross-section in $\Omega_{P_o}^{\times}$ with 
$\Xi_{\cT_{P_o}} \subset \overline{P_o - P_o}$.
\end{lemma}

\begin{lemma}
\label{lemma3}
Let $(F_n)$ be a strong F\o lner sequence in $G$ and suppose $\mu \in \Prob_G(\Omega_{P_o}^{\times})$ is ergodic. Let $g_1,\ldots,g_r \in \Delta$. 
Then, for $\mu$-almost every $P \in \Omega_{P_o}^{\times}$, there exists a 
sub-sequence $(F_{n_j})$ such that
\[
\lim_{j \ra \infty}
\frac{m_G\big(\big(P \cap \big(\bigcap_{k=1}^r (P-g_k)\big)+V_o\big) \cap F_{n_j}\big)}{m_G(F_{n_j})}
= m_G(V_o) \cdot \mu_{\cT_{P_o}}\big(\bigcap_{k=1}^r (-g_k).\cT_{P_o}\big).
\]
\end{lemma}

\begin{proof}
Define the (possibly empty) open set $A \subset \Omega_{P_o}^{\times}$ by
\[
A := \cO_{V_o} \cap \Big( \bigcap_{k=1}^r (-g_k).\cO_{V_o} \Big).
\]
Lemma \ref{Lemma_g1gr} (applied to $B = \cT_{P_o}$, and using that $V_o$ is symmetric) implies that
\[
\mu(A) = m_G(V_o) \cdot \mu_{\cT_{P_o}}\Big(\cT_{P_o} \cap \Big(\bigcap_{k=1}^r (-g_k).\cT_{P_o}\Big)\Big).
\]
Since $\mu$ is ergodic and $(F_n)$ is F\o lner, the mean ergodic theorem tells us that
\[
\lim_{n \ra \infty} \frac{1}{m_G(F_n)} \int_{F_n} \chi_{A}(g.  \, \cdot) \, dm_G(g)
= \mu(A),
\]
in the norm topology on $L^2(\Omega_{P_o}^{\times},\mu)$. We can thus extract  a sub-sequence $(F_{n_j})$ along which the convergence is $\mu$-almost sure.
Hence, for $\mu$-almost every $P \in \Omega_{P_o}^{\times}$,
\[
\lim_{j \ra \infty} \frac{m_G(\{ g \in F_{n_j} \, : \, P - g \in A\})}{m_G(F_{n_j})} 
= \mu(A).
\] 
Note that since $V_o$ is symmetric, 
\begin{align*}
&\{ g \in F_{n_j} \, : \, P - g \in A\} \\[0.2cm]
&=
\{g \in F_{n_j} \, : \, (P - g) \cap V_o \neq \emptyset, \enskip (P-g-g_k) \cap V_o \neq \emptyset, \enskip \textrm{for all $k=1,\ldots,r$}\} \\
&=
(P+V_o) \cap \big( \bigcap_{k=1}^r (P + V_o - g_k) \big) \cap F_{n_j},
\end{align*}
for all $j$. It thus suffices to show that
\[
(P+V_o) \cap \big( \bigcap_{k=1}^r (P + V_o - g_k) \big) 
= \big(P \cap \big( \bigcap_{k=1}^r (P-g_k)\big) + V_o \big).
\]
To do this, first note the $\supseteq$-inclusion is trivial, so it is enough to prove
\begin{equation}
\label{incllr}
(P+V_o) \cap \big( \bigcap_{k=1}^r (P + V_o - g_k) \big) 
\subseteq \big(P \cap \big( \bigcap_{k=1}^r (P-g_k)\big) + V_o \big).
\end{equation}
We may without loss of generality assume that the set on the left-hand side is non-empty, and pick $g$ in this set, allowing us to write
\[
g = p_o + v_o = p_1 + v_1 - g_1 = \ldots = p_r + v_r - g_r,
\]
for some $p_o,\ldots,p_r \in P$ and $v_o,v_1,\ldots,v_r \in V_o$. Then, by Lemma \ref{Lemma_TransSep},
\[
p_k - p_o - g_k = v_o - v_k \in (P-P - \Delta) \cap (V_o-V_o) \subset (\overline{P_o - P_o} - \Delta) \cap W = \{0\},
\]
for all $k$, and thus $v_k = v_o$ and $p_o = p_k - g_k$ for all $k$. In particular, 
\[
g = p_o + v_o \in P \cap \big( \bigcap_{k=1}^r (P-g_k)\big) + V_o.
\]
Since $g$ is arbitrary, this proves \eqref{incllr}, and we are done.
\end{proof}

\end{document}
